\section{Introduction}
Total dissolved inorganic carbon ($\tco$) is a fundamental variable in ocean biogeochemistry, governing air--sea $\mathrm{CO}_2$ flux, ocean acidification, and the biological carbon pump \citep{Sarmiento2006,Takahashi1993}. Quantifying $\tco$ distributions from hydrographic predictors such as salinity ($S$), temperature ($T$), and apparent oxygen utilisation ($\aou$) is vital for gap-filling sparse observational networks, mapping global carbon inventories, and initializing numerical ocean models \citep{Goyet1991,Key2004,Broullon2019}.

Traditional statistical formulations range from simple linear regressions \citep{Millero1998} to complex multi-layer neural networks \citep{Broullon2019} and multi-parameter interpolation tools \citep{Carter2018}. While linear models offer parameter efficiency and mathematical simplicity, they struggle to capture the nonlinear curvature driven by carbonate buffer chemistry and water-mass mixing dynamics. Conversely, non-parametric machine learning models attain high predictive precision but act as opaque black boxes, providing limited insight into the intrinsic degrees of freedom governing empirical carbon variability.

Symbolic regression (SR) has emerged as an automated technique for discovering compact closed-form mathematical equations directly from observational datasets \citep{Cranmer2023}. Rather than treating SR as an automated algorithm for discovering an inviolable "universal law," this study uses symbolic regression as a rigorous probe of empirical low-dimensionality. Specifically, we address the following central question:
\begin{quote}
\emph{How low-dimensional is the empirical $S$--$T$--$\aou$ dependence of ocean total dissolved inorganic carbon, and how much predictive accuracy is retained as structural constraints become stronger?}
\end{quote}

We hypothesize that much of the empirical variability of $\tco$ across the global ocean can be compressed into a low-dimensional nonlinear representation---specifically a five-parameter rank-1 quadratic form---without requiring the full complexity of higher-degree polynomials or deep neural architectures. To test this hypothesis comprehensively, we evaluate candidate expressions across four major ocean basins (Atlantic, Indian, Pacific, and Southern Ocean) using a suite of 48 symbolic regression searches across four distinct mathematical operator spaces and three random seeds.

Importantly, our work explicitly refrains from claiming:
\begin{enumerate}[label=(\roman*)]
\item the discovery of a universal physical ocean-carbon law;
\item a first-principles derivation of aqueous carbonate equilibrium;
\item proof of a single physical mixing axis;
\item that a rank-1 model is universally optimal or invariant under all validation regimes; or
\item that symbolic regression uniquely isolated the single "true" equation governing ocean carbon.
\end{enumerate}

Instead, we systematically evaluate the tradeoff between predictive precision and expression complexity across a broad model hierarchy ranging from mean baselines and linear models up to full quadratic polynomials, rank-2 quadratics, cubic regressions, and Pareto-optimal symbolic candidates. We test model transferability using temporal validation, $5^\circ \times 10^\circ$ spatial block cross-validation, cruise-level cross-validation, and a locked post-2018 external holdout dataset. Finally, we map out specific water-mass and abyssal failure regimes, establishing both the utility and the explicit boundaries of compact ocean carbon parameterizations.
